\documentclass[11pt,a4paper]{article}
\usepackage[utf8]{inputenc}
\usepackage[vietnamese]{babel}
\usepackage{amsmath,amssymb}
\usepackage{geometry}
\geometry{a4paper, total={180mm,260mm}, left=15mm, right=15mm, top=18mm, bottom=18mm}
\usepackage{tikz}
\usetikzlibrary{shapes.geometric, arrows.meta, positioning, calc}
\usepackage{tcolorbox}
\tcbuselibrary{skins,breakable}
\usepackage{xcolor}
\usepackage{multicol}
\usepackage{array}
\usepackage{fontawesome}
\usepackage{fancyhdr}

\pagestyle{fancy}
\fancyhf{}
\renewcommand{\headrulewidth}{0.5pt}
\renewcommand{\footrulewidth}{0.5pt}
\fancyhead[L]{\small\bfseries\color{teal}\faBook\ SỔ TAY KIẾN THỨC TOÁN 3}
\fancyhead[R]{\small\bfseries\color{teal}Lớp Toán Thầy Lĩnh -- loptoanthaylinh.vn}
\fancyfoot[C]{\small\color{gray}-- Sổ tay kiến thức: Bài 33. Nhiệt độ. Đơn vị đo nhiệt độ ($^\circ\text{C}$) --}

\definecolor{primary}{HTML}{005F73}
\definecolor{secondary}{HTML}{0A9396}
\definecolor{accent}{HTML}{EE9B00}
\definecolor{danger}{HTML}{AE2012}
\definecolor{boxbg}{HTML}{F4F6F8}

\begin{document}

% --- BANNER TIÊU ĐỀ ---
\begin{tcolorbox}[
  enhanced,
  colback=white,
  colframe=primary,
  arc=3mm,
  boxrule=1.2pt,
  left=4mm, right=4mm, top=2.5mm, bottom=2.5mm,
  drop shadow=black!10!white
]
  \centering
  {\Large\bfseries\color{primary}\faBookmark\ BÀI 33: NHIỆT ĐỘ. ĐƠN VỊ ĐO NHIỆT ĐỘ ($^\circ\text{C}$)}\\[1.5mm]
  {\small\itshape\color{black!75}Khái niệm nhiệt độ, thang đo độ C, nhiệt kế đo thời tiết, nhiệt kế đo thân nhiệt trên 1 trang A4}
\end{tcolorbox}

\vspace{-2mm}


\noindent
\begin{minipage}[t]{0.49\textwidth}\vspace{0pt}
\begin{tcolorbox}[
  enhanced,
  colback=blue!4!white,
  colframe=blue!75!black,
  arc=2mm, boxrule=0.8pt,
  left=2mm, right=2mm, top=1.8mm, bottom=1.8mm,
  title={\footnotesize\bfseries\color{white}\faKey\ 1. Khái niệm và Đơn vị nhiệt độ}
]
{\footnotesize
\begin{itemize}\setlength{\itemsep}{1.5pt}\setlength{\leftskip}{-2.5mm}
  \item Để đo độ nóng, lạnh của một vật hay thời tiết, người ta dùng đại lượng \textbf{nhiệt độ}.
  \item Đơn vị đo nhiệt độ thường dùng là \textbf{độ C}, kí hiệu là $\mathbf{^\circ\text{C}}$.
  \item Dụng cụ đo nhiệt độ là \textbf{nhiệt kế}.
  \item \textbf{Các mốc nhiệt độ thực tế cần nhớ:}
  \begin{itemize}\setlength{\itemsep}{0.5pt}\setlength{\leftskip}{-2.5mm}
    \item Nước đá đang tan: $\mathbf{0^\circ\text{C}}$.
    \item Nước đang sôi: $\mathbf{100^\circ\text{C}}$.
    \item Thân nhiệt người bình thường: khoảng $\mathbf{37^\circ\text{C}}$.
    \item Phòng điều hòa mát: khoảng $25^\circ\text{C} - 28^\circ\text{C}$.
  \end{itemize}
\end{itemize}
}
\end{tcolorbox}
\end{minipage}
\hfill
\begin{minipage}[t]{0.49\textwidth}\vspace{0pt}
\begin{tcolorbox}[
  enhanced,
  colback=teal!4!white,
  colframe=teal!75!black,
  arc=2mm, boxrule=0.8pt,
  left=2mm, right=2mm, top=1.8mm, bottom=1.8mm,
  title={\footnotesize\bfseries\color{white}\faSliders\ 2. Kĩ năng đọc nhiệt kế}
]
{\footnotesize
\begin{itemize}\setlength{\itemsep}{1.5pt}\setlength{\leftskip}{-2.5mm}
  \item \textbf{Nhiệt kế thủy ngân / rượu:}\\
  Nhìn mực chất lỏng màu đỏ hoặc bạc dừng ở vạch số nào $\implies$ đọc số nhiệt độ tương ứng.
  \item \textbf{Nhiệt kế điện tử:}\\
  Đọc trực tiếp con số hiển thị trên màn hình LCD.
  \item \textbf{So sánh nhiệt độ:}\\
  Số đo lớn hơn $\implies$ nóng hơn / ấm hơn.\\
  Số đo bé hơn $\implies$ lạnh hơn / mát hơn.
\end{itemize}
}
\end{tcolorbox}
\end{minipage}

\vspace{-2.5mm}
% --- KHỐI 3: SƠ ĐỒ TƯ DUY TIKZ ---
\noindent
\begin{tcolorbox}[
  enhanced,
  colback=orange!4!white,
  colframe=orange!80!black,
  arc=2mm, boxrule=0.8pt,
  left=2mm, right=2mm, top=1.5mm, bottom=1.5mm,
  title={\footnotesize\bfseries\color{white}\faSitemap\ SƠ ĐỒ TƯ DUY: THANG NHIỆT ĐỘ VÀ CÁC MỐC THỰC TIỄN}
]
\centering
\begin{tikzpicture}[
  box/.style={draw, rounded corners=2pt, align=center, font=\scriptsize\bfseries, inner sep=2.5pt, minimum height=0.75cm},
  arr/.style={-{Stealth[length=2mm]}, thick}
]
  \node[box, fill=blue!15, draw=blue!80!black, text width=2.6cm] (t1) at (0, 0) {$0^\circ\text{C}$\\Nước đá đang tan\\Rất lạnh / Đóng băng};
  \node[box, fill=teal!15, draw=teal!80!black, text width=2.6cm] (t2) at (3.5, 0) {$20^\circ\text{C} - 28^\circ\text{C}$\\Mát mẻ / Dễ chịu\\Nhiệt độ phòng};
  \node[box, fill=orange!15, draw=orange!80!black, text width=2.6cm] (t3) at (7.0, 0) {$37^\circ\text{C}$\\Thân nhiệt người\\($> 38^\circ\text{C}$ là bị sốt)};
  \node[box, fill=red!15, draw=red!80!black, text width=2.6cm] (t4) at (10.5, 0) {$100^\circ\text{C}$\\Nước đang sôi\\Rất nóng, cẩn thận!};

  \draw[arr, blue!80!black] (t1) -- (t2);
  \draw[arr, teal!80!black] (t2) -- (t3);
  \draw[arr, orange!80!black] (t3) -- (t4);
\end{tikzpicture}
\end{tcolorbox}

\vspace{-2.5mm}
% --- KHỐI 4: VÍ DỤ & MẸO TRÁNH BẪY ---
\noindent
\begin{minipage}[t]{0.49\textwidth}\vspace{0pt}
\begin{tcolorbox}[
  enhanced,
  colback=green!4!white,
  colframe=green!60!black,
  arc=2mm, boxrule=0.8pt,
  left=2mm, right=2mm, top=1.8mm, bottom=1.8mm,
  title={\footnotesize\bfseries\color{white}\faCheckCircle\ Ví dụ thực hành chuẩn mực}
]
{\footnotesize
\textbf{1. Đọc và so sánh nhiệt độ:}\\
Nhiệt độ buổi sáng là $22^\circ\text{C}$, buổi trưa là $31^\circ\text{C}$. Buổi trưa nhiệt độ tăng thêm:\\
$31^\circ\text{C} - 22^\circ\text{C} = \mathbf{9^\circ\text{C}}$.\\
\textbf{2. Nhận biết thực tế:}\\
Khi cặp nhiệt độ thấy chỉ $38^\circ\text{C}5$, bạn Nam đã bị \textbf{sốt} (vì thân nhiệt bình thường là $37^\circ\text{C}$). Cần báo ngay cho cha mẹ hoặc thầy cô.
}
\end{tcolorbox}
\end{minipage}
\hfill
\begin{minipage}[t]{0.49\textwidth}\vspace{0pt}
\begin{tcolorbox}[
  enhanced,
  colback=red!4!white,
  colframe=red!70!black,
  arc=2mm, boxrule=0.8pt,
  left=2mm, right=2mm, top=1.8mm, bottom=1.8mm,
  title={\footnotesize\bfseries\color{white}\faExclamationTriangle\ Mẹo tránh bẫy trắc nghiệm}
]
{\footnotesize
\begin{itemize}\setlength{\itemsep}{1.5pt}\setlength{\leftskip}{-2.5mm}
  \item \textbf{Bẫy quên kí hiệu độ C ($^\circ\text{C}$):}\\
  Khi viết kết quả đo nhiệt độ, bắt buộc phải có kí hiệu $^\circ\text{C}$. Không viết trơ trọi con số $35$ hay $35\text{ C}$.
  \item \textbf{Bẫy đọc vạch nhiệt kế:}\\
  Để ý kỹ mỗi vạch nhỏ tương ứng $1^\circ\text{C}$ hay $2^\circ\text{C}$ để không đọc nhầm số đo.
\end{itemize}
}
\end{tcolorbox}
\end{minipage}


\end{document}
